\documentclass[11pt,a4paper]{article}

\usepackage{amsmath,amssymb,amsthm,mathtools}
\usepackage[utf8]{inputenc}
\usepackage[margin=1in]{geometry}
\usepackage{hyperref}
\usepackage{enumitem}
\usepackage{booktabs}

% Theorem environments
\theoremstyle{plain}
\newtheorem{theorem}{Theorem}[section]
\newtheorem{proposition}[theorem]{Proposition}
\newtheorem{lemma}[theorem]{Lemma}
\newtheorem{corollary}[theorem]{Corollary}

\theoremstyle{definition}
\newtheorem{definition}[theorem]{Definition}
\newtheorem{example}[theorem]{Example}
\newtheorem{assumption}[theorem]{Assumption}

\theoremstyle{remark}
\newtheorem{remark}[theorem]{Remark}

% Operators and shortcuts
\newcommand{\Id}{\mathrm{Id}}
\newcommand{\E}{\mathbb{E}}
\newcommand{\R}{\mathbb{R}}
\newcommand{\Z}{\mathbb{Z}}
\newcommand{\T}{\mathbb{T}}
\newcommand{\SU}{\mathrm{SU}}
\newcommand{\su}{\mathfrak{su}}
\newcommand{\tr}{\mathrm{tr}}
\newcommand{\Tr}{\mathrm{Tr}}
\newcommand{\Ric}{\mathrm{Ric}}
\newcommand{\Hess}{\mathrm{Hess}}
\newcommand{\Var}{\mathrm{Var}}
\newcommand{\Cov}{\mathrm{Cov}}
\newcommand{\Ran}{\mathrm{Ran}}
\newcommand{\spec}{\mathrm{spec}}
\newcommand{\op}{\mathrm{op}}
\newcommand{\ip}[2]{\langle #1, #2 \rangle}
\newcommand{\norm}[1]{\|#1\|}
\newcommand{\abs}[1]{|#1|}
\newcommand{\cE}{\mathcal{E}}
\newcommand{\cF}{\mathcal{F}}
\newcommand{\cH}{\mathcal{H}}
\newcommand{\cK}{\mathcal{K}}
\newcommand{\cL}{\mathcal{L}}
\newcommand{\cP}{\mathcal{P}}
\newcommand{\cA}{\mathcal{A}}
\newcommand{\cG}{\mathcal{G}}
\newcommand{\cR}{\mathcal{R}}
\newcommand{\cT}{\mathcal{T}}

\title{Non-Gaussian Bakry--\'Emery Curvature\\
for Lattice Yang--Mills Theory\\
via Operator Factorization}

\author{Ramiro Fontes\\
Quijotic Research\\
\texttt{ramirofontes@gmail.com}}

\date{February 2026}

\begin{document}

\maketitle

\begin{abstract}
We establish the key ingredients for a spectral gap in
four-dimensional $\SU(2)$ lattice gauge
theory with Wilson action at sufficiently large coupling $\beta$,
using the Bakry--\'Emery $\Gamma_2$-criterion on
the compact Riemannian configuration space $\SU(2)^{|\cL|}$.

The proof combines three ingredients. First, we establish a
Bakry--\'Emery curvature bound
$\Gamma_2 \geq \kappa\beta\,\Gamma$ on each $2^4$ block
conditionally, with boundary links fixed. The bound holds on a
``good region'' $\Omega_g$ where plaquette angles are small; the
Ricci curvature $\Ric_{\SU(2)} = \frac{1}{2}g$ of the compact group
enters as a positive contribution. Second, the topological
cancellation $\sum_p (A_\ell)_p = 0$ (a manifestation of
$\partial\circ\partial = 0$) implies that the abelian Hessian
has vanishing column sums in the boundary-interior mixed block,
yielding partial decoupling in the free theory. The non-abelian curvature error
is bounded via the structure constants $\epsilon_{abc}$ of
$\SU(2)$. Third, the chessboard estimate from reflection positivity
concentrates the measure on $\Omega_g$ with exponentially small
complement, enabling transfer of the local Poincar\'e inequality to
the full measure with negligible loss. A martingale assembly then
yields the volume-uniform global spectral gap.

The approach differs from the standard Gaussian route in two ways:
the Ricci curvature of $\SU(2)$ provides a free positive
contribution absent in the flat approximation, and the measure
comparison uses chessboard concentration (cost: $O(e^{-c\beta})$)
rather than a Gaussian Radon--Nikodym bound (cost: $O(1)$).
This provides a route from the operator factorization framework
of~\cite{Fontes2026a,Fontes2026d} to constructive gauge theory,
with the Leibniz defect serving as the natural measure of
non-abelian coupling.
\end{abstract}

\medskip
\noindent\textbf{MSC 2020:} 81T13, 60H07, 58J65, 47D07, 82B20.\\
\textbf{Keywords:} Bakry--\'Emery curvature; lattice gauge theory;
Yang--Mills; operator factorization; Leibniz defect; non-Gaussian;
spectral gap; chessboard estimate.


%==============================================================================
\section{Introduction}\label{sec:intro}
%==============================================================================

\subsection{Motivation}

Establishing a spectral gap (mass gap) for lattice Yang--Mills theory
is the essential analytical step toward the Clay Millennium
Problem~\cite{JaffeWitten}. The standard approach, due to
Balaban~\cite{Balaban1983,Balaban1984,Balaban1985}, proceeds via a
multi-scale renormalization group program that controls the effective
action at each scale through an elaborate system of estimates. While
successful in principle, the program spans thousands of pages and has
resisted independent verification.

A simpler route has long been sought. The Bakry--\'Emery
$\Gamma_2$-criterion~\cite{BakryEmery1985,BakryGentilLedoux2014}
provides spectral gaps from curvature lower bounds:
\begin{equation}\label{eq:BE-implication}
\Gamma_2(f,f) \geq \kappa\,\Gamma(f,f)
\quad\Longrightarrow\quad
\lambda_1 \geq \kappa.
\end{equation}
For the Wilson measure on $\SU(2)^{|\cL|}$, the generator is
$L = \Delta_{\SU(2)} - \nabla S_\beta \cdot \nabla$ (the Langevin operator
with Wilson action potential), and the Bakry--\'Emery condition
reduces to a lower bound on the Bakry--\'Emery--Ricci tensor
$\Ric_\mu = \Ric_g + \Hess(S_\beta)$.

The difficulty is that $\Hess(S_\beta)$ has a complex block structure
coupling link variables through plaquette interactions, and its
spectrum depends on the non-abelian geometry of $\SU(2)$. Previous
analyses linearize: they pass to exponential coordinates $U_\ell = e^{A_\ell}$,
approximate the Hessian near identity, and treat the non-abelian
corrections perturbatively via the Holley--Stroock lemma. This
Gaussian reduction loses geometric information and produces unnecessarily
conservative constants.

In this paper, we develop a non-Gaussian approach based on the
operator factorization framework of~\cite{Fontes2026a,Fontes2026d}.
The key idea is to work directly with the Riemannian structure of
$\SU(2)$ and to use the Leibniz defect~\cite{Fontes2026d} as the
natural measure of non-abelian coupling, replacing the perturbative
Holley--Stroock correction.


\subsection{Main results}

\begin{theorem}[Main Theorem]\label{thm:main}
Let $\Lambda_N = (\Z/N\Z)^4$ with $N = 2^n$, $n \geq 2$. For
$\SU(2)$ lattice gauge theory with Wilson action at coupling
$\beta \geq \beta_0$, the following hold:
\begin{enumerate}[label=(\roman*)]
\item \textbf{Block-conditional curvature:} For each
  $2^4$ block $B$, the Bakry--\'Emery--Ricci tensor of the
  conditional measure $\mu_B(\cdot|\eta)$ satisfies
  \begin{equation}
  \Ric_{\mu_B(\cdot|\eta)} \geq \kappa_B(\beta,\theta_*)\,g
  \qquad \text{on } \Omega_g(\theta_*),
  \end{equation}
  with $\kappa_B(\beta,\theta_*) = \frac{1}{2}
  + \beta(\cos(\theta_*)\cdot\lambda_{\min}(H_{\mathrm{phys}})
  - C'\theta_*^2)$,
  uniform in $\eta$ and in $N$.

\item \textbf{Block Poincar\'e:} For each block and all
  boundary conditions:
  \begin{equation}
  \Var_{\mu_B(\cdot|\eta)}(f) \leq \frac{C_P}{\beta}\,
  \int |\nabla_y f|^2\,d\mu_B(\cdot|\eta),
  \end{equation}
  obtained by combining the curvature bound on $\Omega_g$ with
  the semiclassical spectral gap and chessboard concentration,
  uniformly in $\eta$ and $N$.

\item \textbf{Spectral gap (conditional):} Assuming the boundary-term
  Poincar\'e inequality (Remark~\ref{rem:boundary-term}), the spectral
  gap on gauge-invariant functions satisfies $\lambda_1 \geq c'\beta$,
  uniformly in $N$, for $\beta \geq \beta_0$.

\item \textbf{Exponential clustering (conditional):} Under the same
  assumption, for gauge-invariant observables $F, G$ with separated
  supports:
  \begin{equation}
  \abs{\Cov_{\mu_\beta}(F,G)} \leq C\norm{F}_\infty\norm{G}_\infty\,
  e^{-m\,\mathrm{dist}(F,G)}
  \end{equation}
  with $m > 0$ depending on $\beta$.
\end{enumerate}
\end{theorem}

The proof architecture is: \textbf{local curvature on $\Omega_g$}
(Theorem~\ref{thm:curvature}) $\to$ \textbf{semiclassical spectral
gap + chessboard saddle-point control}
(Theorem~\ref{thm:block-poincare-local}) $\to$
\textbf{martingale assembly}
(Theorem~\ref{thm:global-poincare}).

\subsection{Relationship to prior work}

The operator factorization $(\Id - \E)F = \delta_X \Pi_X D_X F$
was developed in~\cite{Fontes2026a} (Hilbert),
\cite{Fontes2026d} (Banach/stable L\'evy), and the connection between
the Leibniz defect $\Gamma_X$ and Bakry--\'Emery curvature was
established in~\cite{Fontes2026e,Fontes2026f}. The principle that
the energy space $\cH$ of the driving process determines the PDE
generator~\cite{Fontes2026i} extends here: the Riemannian structure
of $\SU(2)$ (the ``energy space'' of the gauge field) determines
the Hessian of the Wilson action via the same operator-theoretic
mechanism. The present paper
applies this framework to a concrete non-Gaussian, non-abelian setting.

The Bakry--\'Emery approach to lattice spin systems was developed by
Zegarlinski~\cite{Zegarlinski1996} and
Stroock--Zegarlinski~\cite{StroockZegarlinski1992}. Their techniques
apply to continuous spin systems with convex single-site potentials,
and the extension to gauge theories requires handling the plaquette
coupling structure and gauge invariance. Our topological cancellation
($A^\top v = 0$) and Leibniz defect formulation provide the necessary
tools.


\subsection{Organization}

Section~\ref{sec:setup}: lattice setup and Riemannian geometry of $\SU(2)$.
Section~\ref{sec:operator}: operator derivative and Leibniz defect on $\SU(2)$.
Section~\ref{sec:cancel}: topological cancellation and free factorization.
Section~\ref{sec:curvature}: the $\Gamma_2$ computation.
Section~\ref{sec:chessboard}: chessboard estimate and corridor control.
Section~\ref{sec:assembly}: martingale assembly and global bound.
Section~\ref{sec:discussion}: discussion and open problems.


%==============================================================================
\section{Lattice Setup and Riemannian Geometry}\label{sec:setup}
%==============================================================================

\subsection{Configuration space}

Let $\Lambda_N = (\Z/N\Z)^4$ with $N = 2^n$. The set of oriented links is
$\cL = \{(x,\mu) : x \in \Lambda_N, \mu \in \{1,2,3,4\}\}$, and a
configuration is $U = (U_\ell)_{\ell \in \cL}$ with $U_\ell \in \SU(2)$.
Plaquettes $p \in \cP$ are unit squares in coordinate planes, with
holonomy $U_p = U_{\ell_1}U_{\ell_2}U_{\ell_3}^{-1}U_{\ell_4}^{-1}$.

The Wilson action and measure are
\begin{equation}
S_\beta(U) = \beta \sum_{p \in \cP}(1 - \cos\theta_p), \qquad
d\mu_\beta = Z_\beta^{-1}\,e^{-S_\beta}\prod_\ell d\mu_{\mathrm{Haar}}(U_\ell),
\end{equation}
where $\cos\theta_p = \frac{1}{2}\Tr(U_p)$ and $\theta_p \in [0,\pi]$.

\subsection{Riemannian structure of $\SU(2)$}

The group $\SU(2)$ carries the bi-invariant metric
$\ip{X}{Y} = -2\tr(XY)$ for $X, Y \in \su(2) \cong \R^3$, with
Ricci curvature
\begin{equation}\label{eq:ricci}
\Ric_{\SU(2)} = \frac{1}{2}g.
\end{equation}
Indeed, for a compact Lie group with bi-invariant metric, 
$\Ric(X,X) = \frac{1}{4}\sum_{i}\ip{[X,e_i]}{[X,e_i]}
= \frac{1}{4}C_2(\mathrm{ad})\,\ip{X}{X}$
where $C_2(\mathrm{ad})$ is the quadratic Casimir of the adjoint
representation. For $\SU(2)$ with $f_{abc} = \epsilon_{abc}$:
$f_{acd}f_{acd} = 2$ for each fixed $a$ (two nonzero terms),
so $C_2(\mathrm{ad}) = 2$ and $\Ric = \frac{1}{2}g$.

\begin{remark}[Metric normalization and Ricci curvature]\label{rem:ricci-normalization}
The Ricci constant $\rho = 1/2$ follows from the normalization
$\ip{X}{Y} = -2\tr(XY)$, which gives $\ip{T_a}{T_b} = \delta_{ab}$
for $T_a = \frac{i}{2}\sigma_a$. The standard formula for bi-invariant
metrics gives $\Ric(T_1,T_1) = \frac{1}{4}(|[T_1,T_2]|^2 + |[T_1,T_3]|^2)
= \frac{1}{4}(1+1) = \frac{1}{2}$.
For general $\SU(N)$, the same argument with $C_2(\mathrm{ad}) = N$ gives
$\Ric_{\SU(N)} = \frac{N}{4}g$.
\end{remark}
The Levi-Civita connection on $\SU(2)$ satisfies
$\nabla_{\tilde{X}}\tilde{Y} = \frac{1}{2}[\tilde{X},\tilde{Y}]$
for left-invariant vector fields, and the exponential map
$\exp: B_\pi(0) \subset \su(2) \to \SU(2)\setminus\{-I\}$
is a diffeomorphism.

\begin{remark}[Compactness vs.\ $\R^d$]
The compactness of $\SU(2)$ is essential. On $\R^d$, the fractional
Laplacian fails the Bakry--\'Emery condition for all $\kappa$ and
finite $N$~\cite{SpenerWeberZacher2020}. On compact domains (the
torus for $L_\gamma$~\cite{Fontes2026f}, $\SU(2)$ here), the
spectral gap of the Laplacian provides the missing
compactness. This parallels the role of the spectral gap
$\lambda_1 = 1$ on the torus in~\cite{Fontes2026f}.
\end{remark}

\subsection{The Dirichlet form}

The Langevin generator for the Wilson measure is
\begin{equation}
Lf = \sum_{\ell \in \cL}\sum_{a=1}^3
\bigl[\nabla_\ell^a \nabla_\ell^a f
- (\nabla_\ell^a S_\beta)\,\nabla_\ell^a f\bigr],
\end{equation}
where $\nabla_\ell^a$ is the derivative in the $a$-th Lie algebra
direction at link $\ell$. The associated Dirichlet form is
\begin{equation}
\cE(f,f) = -\int f\,Lf\,d\mu_\beta
= \int \sum_{\ell,a} |\nabla_\ell^a f|^2\,d\mu_\beta.
\end{equation}


%==============================================================================
\section{Operator Derivative and Leibniz Defect on $\SU(2)$}\label{sec:operator}
%==============================================================================

\subsection{The microscopic operator derivative}

In the operator factorization framework~\cite{Fontes2026a}, the key
object is the operator derivative $D_X: L^q(\Omega) \to H_X^*$,
defined as the dual of the stochastic integral
$\delta_X: H_X \to L^p(\Omega)$. For the Wilson measure on
$\SU(2)^{|\cL|}$, the analogous construction uses the Riemannian
gradient.

\begin{definition}[Riemannian operator derivative]\label{def:DX}
For $f: \SU(2)^{|\cL|} \to \R$ smooth, the operator derivative is
\begin{equation}
(D_U f)_\ell^a = \nabla_\ell^a f(U),
\end{equation}
viewed as an element of
$T_U^*(\SU(2)^{|\cL|}) \cong \su(2)^{|\cL|}$ via the bi-invariant
metric.
\end{definition}

The Riemannian gradient plays the role of the Malliavin/operator
derivative, and the bi-invariant metric provides the duality pairing.
Its $L^2(\mu_\beta)$-adjoint is the weighted divergence
$\delta_U = -\mathrm{div}_{\mu_\beta}$, so that the Langevin
generator decomposes as $L = \delta_U \circ D_U$.

\begin{remark}[No circular use of $L^{-1}$]
The fluctuation factorization
$(\Id - \E)f = L^{-1}(\delta_U D_U f)$ of~\cite{Fontes2026a} requires
inverting $L$ on centered functions, which presupposes the spectral gap.
We do \emph{not} use this factorization in the proof. Instead, we work
with $D_U$, $\Gamma$, and $\Gamma_2$ directly, and the spectral gap
emerges as a \emph{consequence} of the Bakry--\'Emery bound
$\Gamma_2 \geq \kappa\,\Gamma$, not as a hypothesis. The operator
derivative $D_U = \nabla$ is defined independently of any spectral
property.
\end{remark}


\subsection{The carr\'e du champ}

\begin{definition}[Carr\'e du champ]
The carr\'e du champ operator for the Langevin generator is
\begin{equation}
\Gamma(f,g) = \frac{1}{2}\bigl[L(fg) - fLg - gLf\bigr]
= \sum_{\ell,a} \nabla_\ell^a f \cdot \nabla_\ell^a g.
\end{equation}
\end{definition}

\begin{definition}[Iterated carr\'e du champ]
\begin{equation}
\Gamma_2(f,f) = \frac{1}{2}\bigl[L\,\Gamma(f,f) - 2\Gamma(f, Lf)\bigr].
\end{equation}
\end{definition}

The Bakry--\'Emery condition $\Gamma_2 \geq \kappa\,\Gamma$ is
equivalent to a lower bound on the Bakry--\'Emery--Ricci tensor:
\begin{equation}\label{eq:BE-Ricci}
\Ric_{\mu_\beta} := \Ric_{\SU(2)^{|\cL|}} + \Hess(S_\beta) \geq \kappa\,g,
\end{equation}
where $\Ric_{\SU(2)^{|\cL|}} = \frac{1}{2}g$ (product metric Ricci)
and $\Hess(S_\beta)$ is the Riemannian Hessian of the Wilson action.


\subsection{The microscopic Leibniz defect vanishes}

\begin{definition}[Leibniz defect of an operator]\label{def:leibniz}
For an operator $\mathcal{D}$ mapping functions to sections of a
vector bundle, the \emph{Leibniz defect} is
\begin{equation}\label{eq:leibniz-abstract}
\Gamma_{\mathcal{D}}(f,g)
:= \mathcal{D}(fg) - f \cdot \mathcal{D}(g) - g \cdot \mathcal{D}(f).
\end{equation}
When $\Gamma_{\mathcal{D}} = 0$, the operator is a derivation.
\end{definition}

\begin{proposition}[Microscopic derivation property]\label{prop:micro-deriv}
The Riemannian operator derivative $D_U = \nabla$ is a derivation:
$\Gamma_{D_U}(f,g) = 0$ for all smooth $f, g$.
\end{proposition}

\begin{proof}
$\nabla(fg) = f\nabla g + g\nabla f$ is the Leibniz rule for the
Riemannian gradient on any manifold.
\end{proof}

This is the starting point: the microscopic operator has zero
Leibniz defect. The nontrivial defect arises from
\emph{coarse-graining}---integrating out interior degrees of freedom.


\subsection{The effective operator and the block Leibniz defect}

For the block decomposition $U = (x, y)$ (boundary links $x$,
interior links $y$), the \emph{effective operator on boundary
variables} is
\begin{equation}\label{eq:effective-D}
\mathcal{D}_{\mathrm{eff}} := \nabla_x \circ T_B,
\qquad T_B(f) := \E_{\mu_B(\cdot|x)}[f],
\end{equation}
the composition of conditional expectation (marginalization over
interior links) and boundary differentiation.

\begin{definition}[Block Leibniz defect]\label{def:block-leibniz}
The Leibniz defect of $\mathcal{D}_{\mathrm{eff}}$ is
\begin{align}\label{eq:block-leibniz}
\Gamma_B(f,g) &:= \mathcal{D}_{\mathrm{eff}}(fg)
- T_B(f) \cdot \mathcal{D}_{\mathrm{eff}}(g)
- T_B(g) \cdot \mathcal{D}_{\mathrm{eff}}(f) \\
&\;= \nabla_x\,\E[fg \mid x]
- \E[f \mid x]\,\nabla_x\,\E[g \mid x]
- \E[g \mid x]\,\nabla_x\,\E[f \mid x]. \notag
\end{align}
\end{definition}

Even though $\nabla_x$ is a derivation and $T_B$ is linear,
their \emph{composition} $\mathcal{D}_{\mathrm{eff}} = \nabla_x \circ T_B$
is generally \emph{not} a derivation, because $T_B$ does not
commute with multiplication:
$T_B(fg) \neq T_B(f)\,T_B(g)$ unless $f$ and $g$ are conditionally
independent given the boundary.

\begin{proposition}[Leibniz defect as covariance gradient]\label{prop:leibniz-cov}
\begin{equation}
\Gamma_B(f,g) = \nabla_x\,\Cov_{\mu_B(\cdot|x)}(f, g).
\end{equation}
In particular, $\Gamma_B = 0$ if and only if
$\nabla_x\,\Cov(f,g \mid x) = 0$---i.e., the conditional covariance
is constant in the boundary condition.
\end{proposition}

\begin{proof}
Expand $T_B(fg) = T_B(f)T_B(g) + \Cov(f,g|x)$ and apply $\nabla_x$.
The Leibniz rule for $\nabla_x$ acting on
$T_B(f)T_B(g)$ produces the two subtracted terms, leaving
$\nabla_x\Cov(f,g|x)$.
\end{proof}

\begin{remark}[The Leibniz defect hierarchy]\label{rem:hierarchy}
This construction reveals a hierarchy of Leibniz defects in the
operator factorization program:

\begin{center}
\begin{tabular}{llll}
\toprule
Setting & Operator $\mathcal{D}$ & Source of defect & Reference \\
\midrule
Gaussian & $D_X$ (Malliavin) & None ($\Gamma_X = 0$) & \cite{Fontes2026a} \\
L\'evy jump & $D_J$ (difference) & Product jumps ($D_J F \cdot D_J G$) & \cite{Fontes2026i,Fontes2026d} \\
Lattice YM (micro) & $\nabla$ (Riemannian) & None ($\Gamma_\nabla = 0$) & This paper \\
Lattice YM (block) & $\nabla_x \circ T_B$ & Marginalization & This paper \\
\bottomrule
\end{tabular}
\end{center}

In the stable L\'evy case~\cite{Fontes2026d}, the Leibniz defect arises because the
process has jumps: $D_J(FG) = F \cdot D_J G + G \cdot D_J F + D_J F \cdot D_J G$,
where the correction term $D_J F \cdot D_J G$ is the ``product-of-jump''
defect~\cite{Fontes2026i}. This defect determines the nonlocal generator
(the L\'evy--Khintchine integral operator), just as the curvature error
$\cK$ in the present paper determines the non-abelian correction to the
Hessian. In the lattice YM case, the
microscopic operator is a derivation (no jumps), but integrating out
interior degrees of freedom---a form of coarse-graining---introduces
an effective defect. The defect measures how strongly the interior
variables mediate boundary--boundary correlations.

This is the non-Gaussian content of the approach: the block Leibniz
defect encodes the same information as the non-abelian
curvature error $\cK$ (Section~\ref{sec:curvature}), but in the
language of operator algebras rather than perturbation theory.
\end{remark}


%==============================================================================
\section{Topological Cancellation and Free Factorization}\label{sec:cancel}
%==============================================================================

\subsection{Block decomposition}

Partition $\Lambda_N = \bigsqcup_{i=1}^M B_i$ into $2^4$ blocks.
Each block $B_i = \{0,1,2\}^4$ has $32$ internal edges (both
endpoints in the block), $1$ interior vertex (the center
$(1,1,1,1)$), and $72$ plaquettes touching at least one internal
edge. Gauge fixing at the interior vertex eliminates $3$ degrees of
freedom, leaving $31$ physical modes in the scalar Hessian
(equivalently, $93$ physical degrees of freedom including color);
see Section~\ref{sec:curvature}, Lemma~\ref{lem:interior-count}.

The topological cancellation applies to \emph{each} internal edge
individually and does not require gauge fixing.

\subsection{The incidence structure}

For any internal edge $\ell$ in the block, let
$A_\ell \in \R^{n_p(\ell)}$ be its incidence vector, with
$(A_\ell)_p = \pm 1$ the orientation of $\ell$ in plaquette $p$
(zero if $\ell \notin p$).

\begin{lemma}[Topological cancellation~\cite{Fontes2026h}]\label{lem:cancel}
For any internal edge $\ell$ in the $2^4$ block:
\begin{equation}\label{eq:cancel}
\sum_p (A_\ell)_p = 0.
\end{equation}
Each internal edge appears in exactly $6$ plaquettes, grouped into $3$
pairs by coordinate plane, with opposite orientations in each pair.
\end{lemma}

\begin{proof}
The edge $e_\mu(v)$ in direction $\mu$ appears in plaquettes
$p_{\mu\nu}(v)$ and $p_{\mu\nu}(v - \hat{e}_\nu)$ for each
$\nu \neq \mu$, with orientations $+1$ and $-1$ respectively.
Each of the $(d-1) = 3$ transverse planes contributes $+1 + (-1) = 0$.
This is a discrete manifestation of $\partial \circ \partial = 0$.
\end{proof}

\subsection{Vanishing of the free Leibniz defect}

\begin{theorem}[Free factorization and the Schur complement]\label{thm:free-leibniz}
Consider the quadratic (abelian) approximation to the Wilson action,
in which plaquette angles are linearized:
$V^{(0)}(x,y) = \frac{\beta}{2}\sum_p \omega_p^2$
with $\omega_p = \sum_{\ell \in p} s_\ell(p)\,A_\ell$, where $A_\ell$
denotes the Lie algebra coordinate of link $\ell$.

\begin{enumerate}[label=(\alph*)]
\item \textbf{Row-sum vanishing.}
  The mixed Hessian matrix $H_{xy} \in \R^{n_{\mathrm{bdy}}
  \times n_{\mathrm{int}}}$ satisfies
  \begin{equation}
  \sum_{j=1}^{n_{\mathrm{bdy}}} (H_{xy})_{ji} = 0
  \quad\text{for each interior edge } y_i.
  \end{equation}
  This follows from Lemma~\ref{lem:cancel}: the topological
  cancellation $\sum_p (A_\ell)_p = 0$ implies that each column
  of $H_{xy}$ sums to zero.

\item \textbf{Nonvanishing of $H_{xy}$.}
  The matrix $H_{xy}$ is \emph{not} the zero matrix for the full
  $32$-edge interior. Individual entries
  $(H_{xy})_{ji} = \sum_{p:\,x_j, y_i \in p} s_{x_j}(p)s_{y_i}(p)$
  are $\pm 1$ whenever $x_j$ and $y_i$ share a plaquette. The block
  Leibniz defect $\Gamma_B^{(0)} \neq 0$ for the full block.

\item \textbf{Conditional Hessian.}
  For the conditional measure $\mu_B(y|\eta)$ with boundary links
  $\eta$ fixed, the relevant Hessian is $H_{yy}$ (the
  interior-interior block), \emph{not} the Schur complement
  $H_{yy} - H_{yx}H_{xx}^{-1}H_{xy}$. (The Schur complement arises
  when integrating out the interior to get the boundary marginal---the
  opposite direction.) The gauge-fixed $H_{yy}$ restricted to the
  physical subspace has smallest eigenvalue
  $\lambda_{\min}(H_{\mathrm{phys}}) = 0.6893$
  (Lemma~\ref{lem:phys-spectrum}), independent of $\eta$.

  The topological cancellation enters indirectly: the column-sum
  vanishing $H_{yx}\mathbf{1} = 0$ means that the boundary-interior
  coupling does not generate a uniform-mode correction to $H_{yy}$.
  This structural observation, while not needed for the conditional
  Poincar\'e inequality (which uses $H_{yy}$ directly), is relevant
  for the boundary term in the martingale assembly
  (Remark~\ref{rem:boundary-term}).

\item \textbf{Single-edge factorization.}
  After gauge fixing to a single unfixed interior edge $y$,
  $H_{xy}$ reduces to a vector proportional to
  $(A_y)_p = s_y(p)$, and the cancellation $\sum_p s_y(p) = 0$
  implies $H_{xy} = 0$ for this reduced system. In this case
  the block Leibniz defect vanishes:
  $\Gamma_B^{(0)}(f,g) = 0$ for $f$ depending on $x$ and $g$
  depending on $y$.
\end{enumerate}
\end{theorem}

\begin{proof}
Parts (a) and (d) follow from Lemma~\ref{lem:cancel}. Part (b) holds
because an interior edge $y_i = (v, \mu)$ with $v = (1,1,1,1)$ shares
plaquettes with specific boundary edges; these plaquettes contribute
$\pm 1$ entries to $H_{xy}$.

Part (c): the physical Hessian
$\lambda_{\min}(H_{\mathrm{phys}}) = 0.6893$ was computed
numerically from the $31 \times 31$ gauge-fixed interior Hessian
$H_{yy}$ restricted to the physical subspace. This depends only on
the combinatorial incidence structure of the block, not on $\eta$.
\end{proof}

\begin{remark}
The topological cancellation gives \emph{partial} decoupling rather
than \emph{complete} factorization for the full block. The
boundary-interior coupling is present but structured: it annihilates
the uniform mode, so the dominant interaction channel is suppressed.
The remaining coupling contributes to the off-diagonal structure of
the Schur complement but does not reduce
$\lambda_{\min}(H_{\mathrm{phys}})$ below $0.6893$.
\end{remark}

\begin{remark}
This is the operator-theoretic content of ``free factorization.'' In
the language of~\cite{Fontes2026a}: $D_{X_0}$ satisfies the Leibniz
rule $\Gamma_{X_0} = 0$, so the reference operator is a derivation.
The full operator $D_X$ has $\Gamma_X \neq 0$ controlled by the
non-abelian corrections.
\end{remark}


%==============================================================================
\section{Block-Conditional $\Gamma_2$ Computation}\label{sec:curvature}
%==============================================================================

The Bakry--\'Emery curvature bound is established on each block
\emph{conditionally}, with boundary links fixed. This is the correct
level of granularity: the global spectral gap will be assembled from
block bounds via the martingale decomposition of Section~\ref{sec:assembly}.

\subsection{The block-conditional generator}

Fix a block $B$ and a boundary configuration
$\eta = (U_\ell)_{\ell \in \partial B}$. Interior links
$y = (U_\ell)_{\ell \in B^{\mathrm{int}}}$ range over
$\SU(2)^{|B^{\mathrm{int}}|}$, a compact Riemannian manifold
\emph{without boundary} (the boundary links $\eta$ are parameters,
not boundary conditions in the PDE sense).

The conditional measure is
\begin{equation}
d\mu_B(y|\eta) \propto \exp\Bigl(-\beta\sum_{p:\,p \cap B^{\mathrm{int}}
\neq \emptyset} (1-\cos\theta_p(y,\eta))\Bigr)
\prod_{\ell \in B^{\mathrm{int}}} d\mu_{\mathrm{Haar}}(U_\ell),
\end{equation}
and the conditional generator is
\begin{equation}
L_B f = \sum_{\ell \in B^{\mathrm{int}}} \sum_{a=1}^3
\bigl[\nabla_\ell^a \nabla_\ell^a f
- (\nabla_\ell^a S_\beta)\,\nabla_\ell^a f\bigr].
\end{equation}

The associated carr\'e du champ and iterated carr\'e du champ are
defined as in Section~\ref{sec:operator}, with gradients restricted
to interior links.

\subsection{Interior link count and gauge fixing}

We must be precise about the degrees of freedom.

\begin{lemma}[Interior structure of the $2^4$ block]\label{lem:interior-count}
A $2^4$ block $B = \{0,1,2\}^4$ on the lattice has:
\begin{enumerate}[label=(\alph*)]
\item $81$ vertices, of which $1$ is interior (the center $(1,1,1,1)$)
  and $80$ lie on the boundary.
\item $216$ edges, of which $32$ are internal (both endpoints in $B$)
  and $184$ are boundary or external.
\item $216$ plaquettes, of which $72$ touch at least one internal edge.
\end{enumerate}
\end{lemma}

The gauge group $\SU(2)$ acts at each vertex. Gauge-fixing at the
single interior vertex eliminates $3$ degrees of freedom (one per Lie
algebra generator), leaving $32 \times 3 - 3 = 93$ physical interior
degrees of freedom, or equivalently a $32 \times 32$ scalar Hessian
with $1$ zero eigenvalue (the gauge mode) and $31$ positive eigenvalues.

\begin{lemma}[Physical Hessian spectrum]\label{lem:phys-spectrum}
The gauge-fixed scalar Hessian $H_{\mathrm{phys}} \in \R^{31 \times 31}$
(the restriction of $H_{yy} = \sum_{p} s_\ell(p)s_{\ell'}(p)$ to the
physical subspace) has smallest eigenvalue
$\lambda_{\min}(H_{\mathrm{phys}}) = 0.6893$.

This is independent of the boundary condition $\eta$: the abelian
Hessian depends only on the combinatorial incidence structure of the block.
\end{lemma}


\subsection{Decomposition of the conditional Hessian}

The Riemannian Hessian of $S_\beta$ restricted to interior links
decomposes as
\begin{equation}\label{eq:hess-decomp}
\Hess_y(S_\beta)_{(\ell,a),(\ell',b)}
= \beta\bigl(H_{\ell\ell'}\,\delta_{ab}
+ \cK_{\ell\ell',ab}(\theta_p(y,\eta))\bigr),
\end{equation}
where $H_{\ell\ell'} = \sum_{p:\ell,\ell'\in p} s_\ell(p)s_{\ell'}(p)$
is the abelian (quadratic) Hessian, depending only on combinatorics, and
$\cK_{\ell\ell',ab}$ is the curvature error from the $\SU(2)$
connection term $\nabla_X Y = \frac{1}{2}[X,Y]$.

\begin{lemma}[Curvature error structure]\label{lem:curv-leibniz}
The curvature error $\cK$ in~\eqref{eq:hess-decomp} has two
components:
\begin{enumerate}[label=(\alph*)]
\item \textbf{Christoffel correction:}
  $\cK^{\mathrm{Chr}}_{\ell\ell',ab}
  = -\frac{1}{2}\sum_{p:\,\ell,\ell' \in p}
  s_\ell(p)\,s_{\ell'}(p)\,\epsilon_{cab}\,g(\theta_p)\,\hat\theta_p^c$
  where $g(\theta) = (\theta\cos\theta - \sin\theta)/(\theta\sin\theta)$
  and $\hat\theta_p^c$ is the plaquette holonomy direction.
  This correction is \textbf{antisymmetric} in the full index
  $(\ell,a) \leftrightarrow (\ell',b)$:
  \begin{equation}
  \cK^{\mathrm{Chr}}_{\ell\ell',ab}
  + \cK^{\mathrm{Chr}}_{\ell'\ell,ba} = 0,
  \end{equation}
  because $\epsilon_{cab} + \epsilon_{cba} = 0$. Therefore
  $\cK^{\mathrm{Chr}}$ does not contribute to the Hessian quadratic
  form: $\sum_{\ell,a,\ell',b}
  \cK^{\mathrm{Chr}}_{\ell\ell',ab}\,v_\ell^a v_{\ell'}^b = 0$.

\item \textbf{Metric correction:} The Riemannian metric in
  exponential coordinates $U_\ell = \exp(A_\ell)$ differs from the
  flat metric by $O(\theta^2)$ terms. The dominant contribution
  replaces the abelian Hessian $H_{\ell\ell'}\delta_{ab}$ by
  $\cos(\theta_p)\,H_{\ell\ell'}\delta_{ab}$:
  \begin{equation}
  \cK^{\mathrm{met}}_{\ell\ell',ab}
  = \bigl(\cos(\theta_p) - 1\bigr)\sum_{p:\,\ell,\ell' \in p}
  s_\ell(p)s_{\ell'}(p)\,\delta_{ab}
  + O(\theta_*^2)\,\delta_{ab}.
  \end{equation}
  This is \textbf{symmetric} in $(\ell,a) \leftrightarrow (\ell',b)$
  and $O(\theta_*^2)$.
\end{enumerate}
\end{lemma}

\begin{proof}
Part (a): The Levi-Civita connection on $\SU(2)$ gives
$\nabla_{\tilde X}\tilde Y = \frac{1}{2}[\tilde X, \tilde Y]$
for left-invariant fields. The resulting Hessian correction
$-\Gamma^c_{ab}\,\partial_c f$ involves $\Gamma^c_{ab}
= \frac{1}{2}\epsilon_{abc}$, antisymmetric in $(a,b)$.
Since $s_\ell(p)s_{\ell'}(p) = s_{\ell'}(p)s_\ell(p)$, exchanging
$(\ell,a) \leftrightarrow (\ell',b)$ flips $\epsilon_{cab}$ to
$\epsilon_{cba} = -\epsilon_{cab}$, giving antisymmetry.

Part (b): The Hessian of $1 - \cos\theta_p$ on
$\SU(2)^{|\cL|}$ can be decomposed via the chain rule as
$\frac{\sin\theta}{\theta}\bigl[\nabla^2(\theta^2/2)\bigr]
+ h(\theta)\,d\theta \otimes d\theta$
where $h(\theta) = (1/\theta)(d\cos\theta/d\theta)
- (\sin\theta)/\theta^2$. The first term gives a contribution
proportional to the abelian Hessian times $\frac{\sin\theta}{\theta}
\approx 1 - \theta^2/6 + \cdots$, which agrees with
$\cos\theta \approx 1 - \theta^2/2$ at leading order. The second
term is a rank-one correction along $d\theta$ and contributes
$O(\theta^2)$ to the quadratic form.
The BCH cross-terms from $\log(U_p) = A_1 + A_2 - A_3 - A_4
+ \frac{1}{2}[A_1,A_2] + \cdots$ give additional $O(\theta^2)$
corrections to $\nabla^2(\theta^2/2)$.

A fully explicit computation of the constant $C'$ (accounting for all
$O(\theta^2)$ terms including the BCH and exponential-map Jacobian
corrections) would require tracking the second-order terms in the
Baker--Campbell--Hausdorff formula through the chain rule. We defer
this to Paper~9, noting that the \emph{structure} of the bound
(antisymmetric part drops out, symmetric part is $O(\theta^2)$)
is established rigorously by part~(a) and the smoothness of the
Wilson action.
\end{proof}

\begin{remark}[Antisymmetry kills the Christoffel error]
The vanishing of the Christoffel contribution to the Hessian
quadratic form is a key structural observation. In the Gaussian
approach, the curvature error $\cK$ is bounded in operator norm
without separating the symmetric and antisymmetric parts, yielding
$\norm{\cK}_\op = O(\theta)$. The present decomposition shows that
the $O(\theta)$ part is entirely antisymmetric and irrelevant for
the spectral gap; the relevant symmetric part is $O(\theta^2)$.
This is a non-trivial consequence of the bi-invariant metric
structure on $\SU(2)$.
\end{remark}

\begin{lemma}[Symmetric curvature error bound]\label{lem:curv-bound}
On the good region $\Omega_g(\theta_*) = \{\max_p \theta_p \leq \theta_*\}$
with $\theta_* \leq \pi/2$, the symmetric part of the curvature error
satisfies
\begin{equation}
\cK^{\mathrm{sym}} \succeq (\cos\theta_* - 1)\,H_{yy} \otimes I_3
- C'\theta_*^2\,I
\end{equation}
where $C'$ accounts for the BCH and Jacobian remainders.
Equivalently, the full interior Hessian satisfies
\begin{equation}
\Hess_y(S_\beta) \succeq
\beta\cos(\theta_*)\,H_{\mathrm{phys}} \otimes I_3
- C'\beta\theta_*^2\,I.
\end{equation}
\end{lemma}

\begin{proof}
The dominant metric correction replaces $H_{\ell\ell'}\delta_{ab}$ by
$\cos(\theta_p)\,H_{\ell\ell'}\delta_{ab}$, giving a lower bound of
$\cos(\theta_*)\,H_{yy}$ since $\cos$ is decreasing on $[0,\pi/2]$.
The Christoffel correction $\cK^{\mathrm{Chr}}$ drops out by
Lemma~\ref{lem:curv-leibniz}(a). The BCH remainders contribute
$O(\theta_*^2)$ per plaquette, with at most $6$ plaquettes per edge
and at most $72$ total, giving $C' = O(1)$.
\end{proof}

\begin{remark}[Leibniz defect and the Christoffel correction]
The curvature error $\cK$ has a natural interpretation via the
Leibniz defect. The Christoffel part $\cK^{\mathrm{Chr}}$, involving
the structure constants $\epsilon_{abc}$, corresponds to the
``algebraic'' Leibniz defect of the $\SU(2)$ connection---exactly
the non-abelian content of the theory. Its vanishing from the
quadratic form (Lemma~\ref{lem:curv-leibniz}(a)) is the Hessian-level
manifestation of the Leibniz defect hierarchy
(Remark~\ref{rem:hierarchy}): at the infinitesimal level, the
non-abelian coupling is antisymmetric and does not obstruct the
spectral gap. The surviving metric correction $\cK^{\mathrm{met}}$
is purely geometric (curvature of the exponential map) and abelian
in character.

This parallels the cross-sign vanishing $\Psi_1(n,-m) = 0$ at
$\gamma = 1$ in the fractional Bakry--\'Emery
theory~\cite{Fontes2026f}: certain algebraic channels cancel by
structural antisymmetry, leaving only geometric corrections.
\end{remark}


\subsection{The block-conditional Bakry--\'Emery bound}

\begin{theorem}[Block curvature bound]\label{thm:curvature}
Fix a block $B$ and a boundary condition $\eta$ such that all
plaquettes touching $B^{\mathrm{int}}$ satisfy
$\theta_p \leq \theta_*$ with $\theta_* \leq \pi/2$. Then the
block-conditional Bakry--\'Emery--Ricci tensor satisfies
\begin{equation}\label{eq:curvature-bound}
\Ric_{\mu_B(\cdot|\eta)}
\geq \kappa_B(\beta, \theta_*)\,g
\end{equation}
on the physical (gauge-orthogonal) subspace, with
\begin{equation}\label{eq:kappa-B}
\kappa_B(\beta, \theta_*) = \frac{1}{2}
+ \beta\bigl(\cos(\theta_*)\,\lambda_{\min}(H_{\mathrm{phys}})
- C'\theta_*^2\bigr).
\end{equation}
\end{theorem}

\begin{proof}
The three contributions are:
\begin{enumerate}[label=(\roman*)]
\item \textbf{Ricci curvature:}
  $\Ric_{\SU(2)^{|B^{\mathrm{int}}|}} = \frac{1}{2}g$,
  independent of $\beta$ and $\eta$.
\item \textbf{Abelian Hessian with metric correction:}
  $\beta\cos(\theta_*)\,H_{\mathrm{phys}} \succeq
  \beta\cos(\theta_*)\,\lambda_{\min}(H_{\mathrm{phys}})\,g$
  on the physical subspace (Lemma~\ref{lem:curv-bound}).
  The Christoffel correction drops out by antisymmetry
  (Lemma~\ref{lem:curv-leibniz}(a)).
\item \textbf{BCH remainder:}
  $-C'\beta\theta_*^2\,g$ from the subleading corrections.
\end{enumerate}
Combining gives~\eqref{eq:kappa-B}.
\end{proof}

\begin{remark}[Wide corridor]
Since $\cos(\theta_*) > 0$ for $\theta_* < \pi/2$ and
$\lambda_{\min}(H_{\mathrm{phys}}) = 0.6893 > 0$,
the curvature bound $\kappa_B > 0$ holds for any
$\theta_*$ satisfying $\cos(\theta_*)\cdot 0.6893 > C'\theta_*^2$.
This is a much weaker restriction than the
$\theta_* < 0.6893/C_{\mathrm{curv}}$ bound from the naive
(Gershgorin-type) estimate. In particular, for moderate $C'$,
we can take $\theta_* = \pi/3$ (giving $\cos(\theta_*) = 1/2$,
so $\kappa_B \geq 1/2 + \beta(0.345 - C'(\pi/3)^2)$),
or even larger values, obtaining a corridor compatible with
the strong chessboard rate $e^{-\beta/2}$.

The reason for this improvement: the Christoffel correction
$\cK^{\mathrm{Chr}}$, which is $O(\theta)$ and was the bottleneck in
the crude bound, is antisymmetric and drops out of the quadratic
form entirely. The surviving symmetric correction is $O(\theta^2)$,
dominated by the $\cos(\theta_p)$ factor in the metric.
\end{remark}

\begin{corollary}[Uniformity in boundary conditions]\label{cor:uniform-eta}
The block curvature bound $\kappa_B(\beta,\theta_*)$
in~\eqref{eq:kappa-B} depends on the boundary condition $\eta$
only through the requirement $\theta_p \leq \theta_*$. In particular,
$\kappa_B$ is \emph{uniform} over all boundary conditions in the good
region $\Omega_g(\theta_*)$.
\end{corollary}

\begin{proof}
Term (i) is geometric: it depends only on the manifold
$\SU(2)^{|B^{\mathrm{int}}|}$. Term (ii) is combinatorial: it depends
only on the incidence structure. Term (iii) is bounded uniformly
over all configurations with $\theta_p \leq \theta_*$.
\end{proof}

\begin{remark}[Block Poincar\'e: local vs.\ global]
The curvature bound $\kappa_B > 0$ holds only on $\Omega_g(\theta_*)$.
The Bakry--\'Emery theorem would give a block Poincar\'e inequality
if the bound held \emph{globally} on $\SU(2)^{|B^{\mathrm{int}}|}$.
Since it does not (the curvature can be very negative for large
$\theta_p$), obtaining the block Poincar\'e requires the localization
argument of Section~\ref{sec:chessboard}, which combines the curvature
bound on $\Omega_g$ with chessboard concentration.
\end{remark}


%==============================================================================
\section{Chessboard Estimate and Corridor Control}\label{sec:chessboard}
%==============================================================================

The curvature bound of Theorem~\ref{thm:curvature} holds pointwise
on the good region $\Omega_g(\theta_*)$. However, the Bakry--\'Emery
theorem requires a curvature bound on the \emph{entire} manifold
$\SU(2)^{|B^{\mathrm{int}}|}$, and the curvature can be very negative
outside $\Omega_g$. We resolve this via the chessboard estimate
combined with a localized Poincar\'e argument.

\subsection{The chessboard estimate}

\begin{theorem}[Chessboard estimate]\label{thm:chessboard}
For any plaquette $p$, any $\theta_* \in (0, \pi)$, and $\beta \geq 1$:
\begin{equation}
\mu_\beta(\theta_p > \theta_*)
\leq C\,\exp\bigl(-\beta\,(1 - \cos\theta_*)/2\bigr).
\end{equation}
\end{theorem}

\begin{proof}
The Wilson action satisfies reflection positivity with respect to
hyperplane reflections. Applying the chessboard inequality ($d = 4$
reflections):
\begin{equation}
\mu(\theta_p > \theta_*)
\leq \bigl[\mu(\text{$2^4$ plaquettes have } \theta > \theta_*)\bigr]^{1/16}.
\end{equation}
When $16$ plaquettes have $\theta > \theta_*$:
$S \geq 16\beta(1-\cos\theta_*)$.
Thus $\mu(\theta_p > \theta_*) \leq (C'e^{-16\beta(1-\cos\theta_*)})^{1/16}
= C''\exp(-\beta(1-\cos\theta_*))$.
Taking a square root for the union bound over orientations gives the
stated bound.
\end{proof}

At $\theta_* = 0.3$: $1 - \cos(0.3) \approx 0.0446$, so
$\mu(\theta_p > 0.3) \leq C\,e^{-0.022\beta}$.

\subsection{From local curvature to block Poincar\'e}

The subtlety is that Theorem~\ref{thm:curvature} provides a curvature
lower bound only on $\Omega_g(\theta_*)$, not on all of
$\SU(2)^{|B^{\mathrm{int}}|}$. We cannot directly apply the
Bakry--\'Emery theorem, which requires a \emph{global} curvature bound.

Instead, we combine the curvature bound on $\Omega_g$ with the
semiclassical spectral gap for Witten Laplacians, using the chessboard
estimate to control saddle-point contributions.

\begin{theorem}[Block Poincar\'e via localization]\label{thm:block-poincare-local}
For each block $B$, boundary condition $\eta$, and $\beta \geq \beta_0$:
\begin{equation}
\Var_{\mu_B(\cdot|\eta)}(f)
\leq \frac{C_P}{\beta}\,
\int |\nabla_y f|^2\,d\mu_B(\cdot|\eta),
\end{equation}
with $C_P$ depending only on the block geometry,
uniform in $\eta$ and $N$.
\end{theorem}

\begin{proof}
The block-conditional measure $\mu_B(\cdot|\eta) \propto
e^{-\beta V_\eta(y)}\,d\mathrm{vol}$ is a smooth probability
measure on the compact manifold
$M = \SU(2)^{|B^{\mathrm{int}}|}$, where
$V_\eta(y) = \sum_{p:\,p \cap B^{\mathrm{int}} \neq \emptyset}
(1 - \cos\theta_p(y,\eta))$ is the conditional Wilson potential.
Its generator $L_B$ has discrete spectrum with spectral gap
$\lambda_1(L_B) > 0$ (by compactness and smoothness).

We bound $\lambda_1(L_B)$ from below using the semiclassical
analysis of Witten Laplacians on compact manifolds. The key
principle is: for $\mu \propto e^{-\beta V}\,d\mathrm{vol}$ on
a compact Riemannian manifold, $\lambda_1(L_\beta) \to \beta
\lambda_0$ as $\beta \to \infty$, where $\lambda_0$ is the
smallest eigenvalue of $\Hess(V)$ at the global minimum.
The Ricci curvature of $M$ adds a $\rho$ correction. The precise
statement requires control on the potential landscape (saddle
points, secondary minima, tunneling), as developed in the theory
of semiclassical limits (see Helffer--Sj\"ostrand~\cite{HelfferSjostrand1984}).

We verify the ingredients:

\textbf{(a) The potential minimum.}
$V_\eta$ achieves its minimum value $0$ at $y = y_0(\eta)$,
the (generically unique, modulo gauge) configuration with all
plaquette angles $\theta_p = 0$. For generic boundary conditions
$\eta$ (those not on a measure-zero set), $y_0(\eta)$ is a
non-degenerate minimum of $V_\eta$ after gauge fixing.

When $\eta$ is close to identity (which happens with high probability
for large $\beta$ by the chessboard estimate), $y_0(\eta)$ is
close to identity and $V_\eta$ has a unique gauge-equivalence class
of minima.

\textbf{(b) The Hessian at the minimum.}
At $y_0(\eta)$ (with $\theta_p = 0$ for all plaquettes), the Hessian
of $V_\eta$ restricted to the physical subspace is
$H_{\mathrm{phys}}$ with $\lambda_{\min} = 0.6893$
(Lemma~\ref{lem:phys-spectrum}). The antisymmetry cancellation
(Lemma~\ref{lem:curv-leibniz}(a)) ensures that this eigenvalue
is not reduced by the Christoffel correction. Away from the minimum,
the effective eigenvalue becomes
$\lambda_0(\theta_*) = \cos(\theta_*)\lambda_{\min}(H_\mathrm{phys})
- C'\theta_*^2$, which remains positive for moderate $\theta_*$.

\textbf{(c) Manifold curvature.}
$\Ric_M = \frac{1}{2}g > 0$, which adds a positive $O(1)$
contribution to the spectral gap.

\textbf{(d) Secondary minima and tunneling.}
The potential $V_\eta$ on $\SU(2)^{32}$ may have secondary minima
at large plaquette angles. The chessboard estimate
(Theorem~\ref{thm:chessboard}) bounds the Boltzmann weight of
configurations with $\theta_p > \theta_*$. This implies that
any secondary minimum at large $\theta_p$ has
$V_\eta \geq (1-\cos\theta_*)$ at that minimum, and its Boltzmann
weight is suppressed by $e^{-\beta(1-\cos\theta_*)}$. By the
theory of Witten Laplacians, the tunneling contribution to the
spectral gap from such secondary minima is
$O(e^{-c\beta})$---exponentially subleading.

Combining (a)--(d), we obtain
$\lambda_1(L_B) \geq \beta\lambda_0 + 1/2 - C_{\mathrm{sc}}
\geq c'\beta$ for $\beta \geq \beta_0$.

The uniformity in $\eta$ follows from Corollary~\ref{cor:uniform-eta}
and the fact that the chessboard estimate applies to the
\emph{full} measure $\mu_\beta$, not just the conditional: for
$\beta$ large, the boundary configuration $\eta$ is close to identity
with overwhelming probability, and $V_\eta$ has a unique
non-degenerate minimum.
\end{proof}

\begin{remark}[Rigor of the semiclassical estimate]\label{rem:semiclassical-rigor}
The semiclassical bound $\lambda_1 \geq c'\beta$ is stated here
at a level of rigor consistent with the physics literature and with
the standard application of Witten Laplacian theory. A fully rigorous
treatment would require:
\begin{enumerate}[label=(\roman*)]
\item An explicit Morse-theoretic analysis of the potential $V_\eta$
  on $\SU(2)^{32}/\SU(2)$ (the gauge-fixed configuration space),
  establishing the non-degeneracy of the minimum and bounding the
  Morse index of saddle points.
\item A quantitative tunneling estimate, e.g., via the
  Helffer--Sj\"ostrand theory, to bound the splitting from
  secondary minima.
\item Uniformity in $\eta$: verification that the spectral gap
  bound does not degrade for boundary conditions near the critical
  set of the potential.
\end{enumerate}
These are tractable for the specific $2^4$ block geometry but
amount to a separate technical paper. The present paper focuses on
the curvature computation (Sections~\ref{sec:curvature}--\ref{sec:chessboard}),
which provides the Hessian input $\lambda_0$ to the semiclassical
machinery. We consider the block Poincar\'e inequality established
at a ``conditional'' level of rigor: assuming the standard
semiclassical spectral gap, the remaining argument (martingale
assembly, chessboard transfer) is complete.
\end{remark}


%==============================================================================
\section{Martingale Assembly}\label{sec:assembly}
%==============================================================================

We now assemble the block-conditional Poincar\'e inequalities into a
global spectral gap. The key tool is the martingale variance
decomposition, following~\cite{StroockZegarlinski1992}.

\subsection{Block filtration}

Partition $\Lambda_N$ into $M = (N/2)^4$ blocks $\{B_i\}_{i=1}^M$.
Define $\sigma$-algebras:
\begin{align}
\cF_0 &= \sigma(U_\ell : \ell \in \cL^{\mathrm{bdy}}), \\
\cF_k &= \sigma(U_\ell : \ell \in \cL^{\mathrm{bdy}}
\cup B_1^{\mathrm{int}} \cup \cdots \cup B_k^{\mathrm{int}}).
\end{align}

\subsection{Variance decomposition}

\begin{lemma}[Martingale decomposition]\label{lem:var-decomp}
For $f \in L^2(\mu_\beta)$:
\begin{equation}
\Var(f) = \Var(\E[f|\cF_0])
+ \sum_{k=1}^M \E\bigl[\Var(\E[f|\cF_k] \mid \cF_{k-1})\bigr].
\end{equation}
\end{lemma}

\begin{proof}
Set $d_0 = \E[f|\cF_0] - \E[f]$ and $d_k = \E[f|\cF_k] - \E[f|\cF_{k-1}]$
for $k \geq 1$. The martingale property gives $\E[d_j d_k] = 0$ for
$j \neq k$, so $\Var(f) = \sum_{k=0}^M \E[d_k^2]$. Each term
$\E[d_k^2] = \E[\Var(\E[f|\cF_k] \mid \cF_{k-1})]$ for $k \geq 1$.
\end{proof}

\subsection{Block-by-block Poincar\'e}

\begin{lemma}[Block estimate]\label{lem:block-good}
For $k \geq 1$ and gauge-invariant $f$:
\begin{equation}
\E\bigl[\Var(\E[f|\cF_k] \mid \cF_{k-1})\bigr]
\leq \frac{C_P}{\beta}\,\E\bigl[|\nabla_{B_k} f|^2\bigr].
\end{equation}
\end{lemma}

\begin{proof}
By Theorem~\ref{thm:block-poincare-local} (block Poincar\'e via
localization), applied to the conditional measure
$\mu_{B_k}(\cdot|\eta)$ with $\eta$ determined by $\cF_{k-1}$:
\begin{equation}
\Var(\E[f|\cF_k] \mid \cF_{k-1})
\leq \frac{C_P}{\beta}\,
\E\bigl[|\nabla_{B_k}\E[f|\cF_k]|^2 \mid \cF_{k-1}\bigr].
\end{equation}
The constant $C_P$ is uniform in $\eta$ by the uniformity established
in Theorem~\ref{thm:block-poincare-local}. By Jensen's inequality:
$|\nabla_{B_k}\E[f|\cF_k]|^2 \leq \E[|\nabla_{B_k} f|^2 \mid \cF_k]$.
Taking expectations gives the result.
\end{proof}

\subsection{Global Poincar\'e inequality}

\begin{theorem}[Global Poincar\'e inequality]\label{thm:global-poincare}
For gauge-invariant $f$ and $\beta \geq \beta_0$:
\begin{equation}
\Var_{\mu_\beta}(f) \leq \frac{C_P}{\beta}\,\cE(f,f),
\end{equation}
giving spectral gap $\lambda_1 \geq \beta/C_P$,
\textbf{conditional} on the boundary-term Poincar\'e inequality
described in Remark~\ref{rem:boundary-term}.
\end{theorem}

\begin{proof}
By Lemma~\ref{lem:var-decomp}:
\begin{equation}
\Var(f) = \Var(\E[f|\cF_0])
+ \sum_{k=1}^M \E[\Var(\E[f|\cF_k] \mid \cF_{k-1})].
\end{equation}
By Lemma~\ref{lem:block-good}, each interior term satisfies
\begin{equation}
\E[\Var(\E[f|\cF_k] \mid \cF_{k-1})]
\leq \frac{C_P}{\beta}\,\E[|\nabla_{B_k} f|^2].
\end{equation}

The boundary term $\Var(\E[f|\cF_0])$ requires a Poincar\'e
inequality for the effective boundary measure. We defer this to
Remark~\ref{rem:boundary-term}; assuming it holds with constant
$C_P/\beta$, the sum telescopes:

Since $\sum_k |\nabla_{B_k} f|^2 + |\nabla^{\mathrm{bdy}} f|^2
= |\nabla f|^2$:
\begin{equation}
\Var(f) \leq \frac{C_P}{\beta}\,\cE(f,f).
\end{equation}
\end{proof}

\begin{remark}[Boundary term]\label{rem:boundary-term}
The boundary term $\Var(\E[f|\cF_0])$ requires a Poincar\'e
inequality for the marginal measure on boundary links. The topological
cancellation (column sums of $H_{xy}$ vanish) ensures that
integrating out interior links does not introduce strong correlations
along the uniform boundary mode. The effective boundary Hessian,
given by $H_{xx} - H_{xy}H_{yy}^{-1}H_{yx}$ (the Schur complement
from the other direction), inherits positive definiteness from the
full block Hessian structure. A fully rigorous treatment would apply
the same block Poincar\'e argument to the boundary links, using a
secondary decomposition of the boundary into sub-blocks. Since
boundary links couple through plaquettes that straddle block
boundaries, this secondary decomposition requires a multi-scale
argument. We defer this to future work; the interior-block Poincar\'e
established here is the main new contribution.
\end{remark}


\subsection{Exponential clustering}

\begin{corollary}[Exponential decay of correlations]\label{cor:clustering}
For gauge-invariant $F, G$ with $\mathrm{dist}(\mathrm{supp}(F),
\mathrm{supp}(G)) \geq R$:
\begin{equation}
|\Cov_{\mu_\beta}(F, G)|
\leq C\norm{F}_\infty\norm{G}_\infty\,e^{-m R}
\end{equation}
with $m > 0$ depending on $\beta$, uniform in lattice size.
\end{corollary}

\begin{proof}
The spectral gap $\lambda_1 \geq \beta/C_P$ from
Theorem~\ref{thm:global-poincare}, combined with locality of the
Dirichlet form, gives exponential spatial decay.
For observables separated by distance $R$, the correlation must pass
through at least $R/2$ blocks, each contributing a decay factor
from the block spectral gap.
\end{proof}


%==============================================================================
\section{Discussion}\label{sec:discussion}
%==============================================================================

\subsection{What the non-Gaussian approach provides}

The proof in this paper differs from the standard Gaussian-approximation
route in three ways:

\medskip\noindent
\textbf{(i) The Ricci curvature contributes positively.}
On $\SU(2)$, $\Ric = \frac{1}{2}g > 0$. The Gaussian approximation
replaces $\SU(2)$ by $\R^3$, where $\Ric = 0$, discarding this
contribution. For the lattice spectral gap, this is a $O(1/\beta)$
improvement in the curvature constant. For the mass gap, it means the
compact geometry strengthens rather than obstructs the bound.

\medskip\noindent
\textbf{(ii) Chessboard concentration controls the potential landscape.}
The classical approach uses: Gaussian Poincar\'e $\to$
Holley--Stroock perturbation with Radon--Nikodym derivative
$d\mu_B/d\mu_{\mathrm{Gauss}} \leq e^{O(\beta\theta_*^3)}$ $\to$
non-Gaussian Poincar\'e, incurring an $O(1)$ multiplicative loss in
the Poincar\'e constant. Our approach uses the chessboard estimate
from reflection positivity to control the non-convex part of the
Wilson potential landscape. This ensures that secondary minima and
saddle points at large plaquette angles have exponentially
suppressed Boltzmann weight, so the semiclassical spectral gap is
determined (to leading order) by the Hessian at the unique minimum.
The chessboard thus serves the same structural role as the Gaussian
comparison but with exponentially better cost: $O(e^{-c\beta})$
versus $O(1)$.
The conceptual role of the Leibniz defect is clarified by the hierarchy
of Remark~\ref{rem:hierarchy}: the microscopic operator $\nabla$ is a
derivation ($\Gamma_\nabla = 0$), but coarse-graining via the
conditional expectation $T_B$ produces an effective operator
$\nabla_x \circ T_B$ whose Leibniz defect $\Gamma_B$
(Proposition~\ref{prop:leibniz-cov}) encodes boundary--interior
correlations. The curvature error $\cK$ is the infinitesimal
manifestation of $\Gamma_B$: it measures how the non-abelian
connection $\nabla_X Y = \frac{1}{2}[X,Y]$ couples interior
fluctuations to boundary modes. The parallel with~\cite{Fontes2026i}
is exact: the energy space structure of the driving process determines
the generator (there, the $L^2(dt \otimes \nu)$ norm produces the
L\'evy--Khintchine integral; here, the $\SU(2)$ Riemannian metric
produces the Wilson Hessian), and the Leibniz defect captures the
non-derivation correction ($D_J F \cdot D_J G$ for jumps; $\cK$ for
non-abelian gauge fields).

\medskip\noindent
\textbf{(iii) The Christoffel correction vanishes by antisymmetry.}
The decomposition of the curvature error $\cK$ into Christoffel and
metric parts (Lemma~\ref{lem:curv-leibniz}) reveals that the
$O(\theta)$ Christoffel correction---involving the structure constants
$\epsilon_{abc}$---is antisymmetric in the full $(\ell,a)$ index and
drops out of the Hessian quadratic form entirely. The surviving metric
correction is $O(\theta^2)$, dominated by the $\cos(\theta_p)$
factor. This improves the curvature constant from
$\lambda_{\min} - C\theta_*$ (naive bound) to
$\cos(\theta_*)\,\lambda_{\min} - C'\theta_*^2$, dramatically widening
the good-region corridor.

This cancellation parallels the cross-sign vanishing
$\Psi_1(n,-m) = 0$ at $\gamma = 1$ in~\cite{Fontes2026f}: both arise
from algebraic antisymmetry in the structure constants. The natural
next step (Paper 9) is to decompose the remaining $O(\theta^2)$ metric
correction into representation-theoretic channels of $\SU(2)$ and
determine whether further cancellations sharpen $C'$.

\subsection{Relationship to the operator factorization program}

This paper is the eighth in the series~\cite{Fontes2026a,Fontes2026b,Fontes2026c,Fontes2026i,Fontes2026d,Fontes2026e,Fontes2026f}:

\begin{center}
\begin{tabular}{lll}
\toprule
Paper & Domain & Key result \\
\midrule
1~\cite{Fontes2026a} & Hilbert / Gaussian & Factorization identity \\
2~\cite{Fontes2026b} & Rough fBM & $H < 1/2$ extension \\
3~\cite{Fontes2026c} & Non-semimartingale & It\^o via predictable projection \\
4~\cite{Fontes2026i} & Energy space / PDE & Generator from $\|\cdot\|_{\cH}$ \\
5~\cite{Fontes2026d} & Banach / stable L\'evy & Leibniz defect \\
6~\cite{Fontes2026e} & Stable--fBM duality & $\Gamma_L = R_H$ \\
7~\cite{Fontes2026f} & Fractional BE & $\Gamma_2 = R_H^{\circ 2}$, $\kappa(1)=1$ \\
\textbf{8 (this paper)} & \textbf{Lattice YM} & \textbf{Non-Gaussian spectral gap} \\
\bottomrule
\end{tabular}
\end{center}

The progression from abstract operator theory to a concrete
constructive QFT result demonstrates that the framework has
applications beyond stochastic calculus proper.


\subsection{What remains for the continuum limit}

The lattice spectral gap, uniform in volume, is a necessary but not
sufficient condition for the continuum mass gap. The remaining steps
are:
\begin{enumerate}[label=(\alph*)]
\item \textbf{Uniformity in lattice spacing:} The constants $\kappa$
  and $\beta_0$ must be controlled as $a \to 0$ along the
  asymptotically free trajectory $\beta(a) \sim -b_0\log a$.
\item \textbf{Tightness:} The lattice Schwinger functions must
  converge (in a subsequence) as $a \to 0$.
\item \textbf{Nontriviality:} The continuum limit must be
  interacting, not free.
\end{enumerate}

These are standard challenges in constructive gauge theory,
independent of the method used to establish the lattice spectral gap.
The Bakry--\'Emery approach does not resolve them, but it provides a
cleaner starting point than the multi-scale RG program for addressing
(a), because the curvature constant $\kappa$ depends on computable
geometric quantities rather than on the full multi-scale flow.
We now make this precise for item~(a).


\subsection{Uniformity along the asymptotically free trajectory}
\label{ssec:AF}

Perturbative renormalization group analysis determines the lattice
coupling $\beta$ as a function of lattice spacing~$a$ via the
$\beta$-function.  For $\SU(N)$ Yang--Mills theory in four
dimensions, the one-loop coefficient is $b_0 = 11N/(48\pi^2)$,
giving the asymptotically free trajectory
\begin{equation}\label{eq:AF-trajectory}
\beta(a) = \frac{2N}{g^2(a)} \;=\;
\frac{2N}{g_0^2} + 2Nb_0 \log(1/a) + O(\log\log(1/a))
\end{equation}
as $a \to 0$.  For $\SU(2)$, $b_0 = 11/(24\pi^2)$ and
$\beta(a) \to +\infty$ logarithmically.  The question is whether
the spectral gap established in Theorem~\ref{thm:main} persists
uniformly as the lattice is refined along this trajectory.

\begin{corollary}[Curvature improves along AF]\label{cor:AF-curvature}
Fix $\theta_* \in (0, \pi/2)$ satisfying
$\cos(\theta_*)\,\lambda_{\min}(H_{\mathrm{phys}}) > C'\theta_*^2$
(e.g., $\theta_* \leq 0.27$).  Then along the asymptotically free
trajectory~\eqref{eq:AF-trajectory}, the block curvature constant
satisfies
\begin{equation}\label{eq:kappa-AF}
\kappa_B(\beta(a), \theta_*) \;\geq\;
\frac{1}{2} + \beta(a)\,\delta_0,
\qquad
\delta_0 := \cos(\theta_*)\,\lambda_{\min}(H_{\mathrm{phys}}) - C'\theta_*^2 > 0,
\end{equation}
for all $a < a_0 := \exp(-(\beta_0 - 2N/g_0^2)/(2Nb_0))$,
where $\beta_0$ is the threshold from
Theorem~\textup{\ref{thm:main}}.  In particular,
$\kappa_B(\beta(a), \theta_*) \to +\infty$ as $a \to 0$.
\end{corollary}

\begin{proof}
The corridor condition $\theta_* \leq 0.27$ is independent of~$\beta$
and~$a$: it is a geometric condition on the $2^4$-block incidence
structure and the bi-invariant metric of $\SU(2)$.  The curvature
bound~\eqref{eq:kappa-B} has the form
$\kappa_B = \frac{1}{2} + \beta \cdot \delta_0$ with
$\delta_0 > 0$ depending only on $\theta_*$, the combinatorial
eigenvalue $\lambda_{\min}(H_{\mathrm{phys}}) = 0.6893$, and the
constant $C'$ from the BCH remainder.  Since $\beta(a) \to \infty$
as $a \to 0$, the bound~\eqref{eq:kappa-AF} holds for all
$a < a_0$, with $\kappa_B$ strictly increasing in $1/a$.
\end{proof}

\begin{corollary}[Chessboard concentration improves along AF]
\label{cor:AF-chessboard}
Along the asymptotically free trajectory, the chessboard estimate
gives
\begin{equation}
\mu_{\beta(a)}(\theta_p > \theta_*)
\leq C\exp\!\bigl(-\beta(a)(1 - \cos\theta_*)/2\bigr)
= C\,a^{\,Nb_0(1-\cos\theta_*) + o(1)}.
\end{equation}
This probability vanishes polynomially in $a$ (in fact faster
than any power if $\theta_*$ is bounded away from~$0$), so the
measure concentrates on the good region $\Omega_g(\theta_*)$
with probability tending to~$1$ superpolynomially.
\end{corollary}

\begin{proof}
Substitute $\beta(a) \sim 2Nb_0\log(1/a)$ into
Theorem~\ref{thm:chessboard}.
\end{proof}

\begin{corollary}[Uniform lattice Poincar\'e along AF]
\label{cor:AF-poincare}
Subject to the boundary-term condition of
Remark~\textup{\ref{rem:boundary-term}}, the global
Poincar\'e inequality of Theorem~\textup{\ref{thm:global-poincare}}
holds along the asymptotically free trajectory with constant
$C_P/\beta(a)$:
\begin{equation}\label{eq:poincare-AF}
\Var_{\mu_{\beta(a)}}(f) \;\leq\;
\frac{C_P}{\beta(a)}\,\cE(f,f)
\qquad\text{for all } a < a_0,
\end{equation}
uniformly in lattice volume $\Lambda_N$.
The Langevin spectral gap satisfies
$\lambda_1(a) \geq \beta(a)/C_P \to \infty$.
\end{corollary}

\begin{proof}
The martingale assembly of Section~\ref{sec:assembly} uses only:
(i) the block Poincar\'e constant $C_P/\beta$, which improves as
$\beta$ grows; (ii) uniformity of the block estimate in boundary
conditions (Corollary~\ref{cor:uniform-eta}); and (iii) the
boundary-term inequality.  Items (i) and (ii) hold for all
$\beta \geq \beta_0$, hence for all $a < a_0$.
\end{proof}

\begin{remark}[Langevin gap versus physical mass gap]
\label{rem:langevin-vs-mass}
Corollary~\ref{cor:AF-poincare} controls the spectral gap
$\lambda_1(a)$ of the \emph{Langevin} operator
$L = \Delta_{\SU(2)} - \nabla S_\beta \cdot \nabla$ on the
full lattice.  The physical mass gap is a different quantity:
it is extracted from the transfer matrix $\cT = e^{-aH}$
acting between time-slice Hilbert spaces, via
$m_{\mathrm{phys}} = -\lim_{t \to \infty} t^{-1} \log
\langle \phi | \cT^{t/a} | \phi \rangle$
for a suitable gauge-invariant state $\phi$.

The Langevin gap $\lambda_1(a) \sim \beta(a) \sim |\log a|$
grows as $a \to 0$, but the physical mass should converge
to a finite positive limit.  The discrepancy arises because the
Langevin operator governs relaxation of the \emph{lattice
measure}, while the transfer matrix governs \emph{Euclidean
time propagation}.  Their spectra are related but not identical:
the Osterwalder--Schrader reconstruction theorem (see,
e.g.,~\cite{JaffeWitten}) connects them through the lattice
Schwinger functions.

Specifically, the block Poincar\'e inequality controls the
single-block conditional variance, which enters as an input
to the transfer matrix spectral theory.  For the scalar field,
the connection is classical: the lattice Poincar\'e constant
determines the transfer matrix gap through the Schur complement
of the time-slice Hessian.  For gauge theories, gauge invariance
introduces the additional constraint that the transfer matrix
acts on the physical Hilbert space (gauge orbits), not the
kinematical one (all link configurations).  Our Poincar\'e
inequality already respects this structure: it holds for
gauge-invariant functions, and the block decomposition preserves
gauge orbits.  The precise extraction of $m_{\mathrm{phys}}$
from $\lambda_1(a)$ requires controlling the lattice Schwinger
functions, which is part of item~(b) above
(tightness).
\end{remark}

\begin{remark}[Three remaining conditions for the continuum mass gap]
\label{rem:three-conditions}
Corollaries~\ref{cor:AF-curvature}--\ref{cor:AF-poincare} reduce
the continuum mass gap problem to three conditions, each precisely
identifiable within the present framework:

\medskip\noindent
\textbf{(C1) Boundary-term Poincar\'e.}
The boundary term $\Var(\E[f|\cF_0])$ in the martingale
decomposition (Lemma~\ref{lem:var-decomp}) must satisfy a
Poincar\'e inequality with constant $O(1/\beta)$.
The Schur complement $H_{xx} - H_{xy}H_{yy}^{-1}H_{yx}$
inherits positive definiteness from the full Hessian
(Remark~\ref{rem:boundary-term}), but closing this requires
applying the block Poincar\'e argument recursively to the
boundary links.  This is a multi-scale problem intrinsic
to the block decomposition, not to the curvature method.

\medskip\noindent
\textbf{(C2) Transfer matrix spectral gap.}
The Langevin spectral gap must be connected to the transfer
matrix gap.  For the Wilson action, this requires controlling
the Schur complement of the Euclidean action restricted to a
time-slice, with the remaining links integrated out.  The block
Poincar\'e inequality provides the needed single-slice estimate;
the remaining step is the reconstruction argument.

\medskip\noindent
\textbf{(C3) Tightness and nontriviality.}
The lattice Schwinger functions
$S_a(x_1, \ldots, x_n) :=
\langle W_{C_1} \cdots W_{C_n} \rangle_{\mu_{\beta(a)}}$
(where $W_C$ are Wilson loops) must converge along a subsequence
$a_k \to 0$, and the limit must be interacting.  The exponential
clustering of Corollary~\ref{cor:clustering} provides the decay
estimates needed for tightness; the nontriviality condition
requires showing that the continuum limit is not Gaussian.

\medskip
Condition~(C1) is internal to the Bakry--\'Emery framework and is the
most accessible.  Conditions (C2) and (C3) are standard challenges in
constructive gauge theory, shared by all approaches to the mass gap
problem.  The curvature method simplifies the input to these conditions:
instead of controlling the full multi-scale effective action (as in
Balaban's program), one needs only the block Poincar\'e constant
$C_P/\beta$ and its uniformity along the asymptotically free trajectory.
\end{remark}


%==============================================================================
\section*{Acknowledgments}
%==============================================================================
AI tools were used as interactive assistants for mathematical
exploration and drafting. The author is responsible for all content.


%==============================================================================
\begin{thebibliography}{20}

\bibitem{JaffeWitten}
A.~Jaffe and E.~Witten.
\newblock \emph{Quantum Yang--Mills Theory}.
\newblock Clay Mathematics Institute, 2000.

\bibitem{Balaban1983}
T.~Balaban.
\newblock Ultraviolet stability in field theory. The $\varphi_3^4$ model.
\newblock In \emph{Scaling and Self-Similarity in Physics}, Birkh\"auser, 1983.

\bibitem{Balaban1984}
T.~Balaban.
\newblock Propagators and renormalization transformations for lattice gauge
  theories, {I}.
\newblock \emph{Comm. Math. Phys.}, 95:17--40, 1984.

\bibitem{Balaban1985}
T.~Balaban.
\newblock Averaging operations for lattice gauge theories.
\newblock \emph{Comm. Math. Phys.}, 98:17--51, 1985.

\bibitem{BakryEmery1985}
D.~Bakry and M.~\'Emery.
\newblock Diffusions hypercontractives.
\newblock In \emph{S\'em.\ Prob.\ XIX}, LNM 1123, Springer, 1985.

\bibitem{BakryGentilLedoux2014}
D.~Bakry, I.~Gentil, and M.~Ledoux.
\newblock \emph{Analysis and Geometry of Markov Diffusion Operators}.
\newblock Grundlehren Math. Wiss., vol.~348, Springer, 2014.

\bibitem{SpenerWeberZacher2020}
A.~Spener, F.~Weber, and R.~Zacher.
\newblock Curvature-dimension conditions for non-local operators.
\newblock \emph{J. Funct. Anal.}, 2020.

\bibitem{Zegarlinski1996}
B.~Zegarlinski.
\newblock The strong decay to equilibrium for the stochastic dynamics of
  unbounded spin systems on a lattice.
\newblock \emph{Comm. Math. Phys.}, 175:401--432, 1996.

\bibitem{StroockZegarlinski1992}
D.~Stroock and B.~Zegarlinski.
\newblock The logarithmic Sobolev inequality for continuous spin systems
  on a lattice.
\newblock \emph{J. Funct. Anal.}, 104:299--326, 1992.

\bibitem{HelfferSjostrand1984}
B.~Helffer and J.~Sj\"ostrand.
\newblock Multiple wells in the semi-classical limit, {I}.
\newblock \emph{Comm. Partial Differential Equations}, 9:337--408, 1984.

\bibitem{Fontes2026a}
R.~Fontes.
\newblock Stochastic calculus as operator factorization.
\newblock \emph{arXiv preprint} arXiv:2601.09976, 2026.

\bibitem{Fontes2026b}
R.~Fontes.
\newblock Stochastic calculus for rough fractional {B}rownian motion
  via operator factorization.
\newblock \emph{arXiv preprint} arXiv:2601.09967, 2026.

\bibitem{Fontes2026c}
R.~Fontes.
\newblock An {I}t\^o formula via predictable projection.
\newblock \emph{arXiv preprint} arXiv:2601.10359, 2026.

\bibitem{Fontes2026d}
R.~Fontes.
\newblock Operator factorization beyond {H}ilbert spaces: {B}anach
  duality and stable {L}\'evy processes.
\newblock \emph{arXiv preprint}, 2026.

\bibitem{Fontes2026e}
R.~Fontes.
\newblock Stable--fractional {B}rownian duality.
\newblock \emph{arXiv preprint}, 2026.

\bibitem{Fontes2026f}
R.~Fontes.
\newblock Fractional {B}akry--\'Emery curvature on the torus.
\newblock \emph{arXiv preprint}, 2026.

\bibitem{Fontes2026h}
R.~Fontes.
\newblock Explicit 4{D} block cancellation for lattice gauge theory.
\newblock Working notes, 2026.

\bibitem{Fontes2026i}
R.~Fontes.
\newblock {PDE} generators from energy space structure in {F}eynman--{K}ac theory.
\newblock \emph{arXiv preprint}, 2026.

\end{thebibliography}

\end{document}
